\subsection{DIRECT SUM OF MODULES}

Let \((\mathbf{E}_t)_{t \in I}\) be a family of A-modules and \(\mathbf{F} = \prod_{i \in I} \mathbf{E}_i\) their product. The set \(\mathbf{E}\) of \(x \in F\) such that \(\operatorname{pr}_t x = 0\) except for a finite number of indices is obviously a submodule of \(\mathbf{F}\) , called the external direct sum (or simply direct sum) of the family of modules \((\mathbf{E}_t)\) and denoted by \(\bigoplus_{i \in I} \mathbf{E}_i\) (I, §4, no. 9). When I is finite, then \(\bigoplus_{i \in I} \mathbf{E}_i = \prod_{i \in I} \mathbf{E}_t\) ; if \(\mathbf{I} = [p, q]\) (interval of \(\mathbf{Z}\) ), we also write

\[
\bigoplus_ {i \in I} \mathbf {E} _ {i} = \mathbf {E} _ {p} \oplus \mathbf {E} _ {p + 1} \oplus \dots \oplus \mathbf {E} _ {q}.
\]

For all \(\kappa \in I\) , let \(j_{\kappa}\) be the mapping \(\mathbf{E}_{\kappa} \to F\) which associates with each \(x_{\kappa} \in \mathbf{E}_{\kappa}\) the element of \(F\) such that \(\operatorname{pr}_{\iota}(j_{\kappa}(x_{\kappa})) = 0\) for \(\iota \neq \kappa\) and \(\operatorname{pr}_{\kappa}(j_{\kappa}(x_{\kappa})) = x_{\kappa}\) ; it is immediate that \(j_{\kappa}\) is an injective linear mapping of \(\mathbf{E}_{\kappa}\) into the direct sum \(E\) of the \(\mathbf{E}_{\iota}\) , which we shall call the canonical injection; the submodule \(j_{\kappa}(\mathbf{E}_{\kappa})\) of \(E\) , isomorphic to \(\mathbf{E}_{\kappa}\) , is called the component submodule of \(E\) of index \(\kappa\) . It is often identified with \(\mathbf{E}_{\kappa}\) by means of \(j_{\kappa}\) .

For all \(x \in \mathbf{E} = \bigoplus_{i \in I} \mathbf{E}_i\) , we have therefore

\[
x = \sum_ {i \in I} j _ {i} (\mathrm{pr} _ {i} x).\tag{20}
\]

PROPOSITION 6. Let \((\mathbf{E}_t)_{t \in I}\) be a family of A-modules, M an A-module and, for all \(t \in I\) , let \(f_t: \mathbf{E}_t \to \mathbf{M}\) be a linear mapping. Then there exists one and only one linear mapping \(g: \bigoplus_{t \in I} \mathbf{E}_t \to \mathbf{M}\) such that, for all \(t \in I\) :

\[
g \circ j _ {i} = f _ {i}.\tag{21}
\]

By virtue of (20), if \(g\) exists, then necessarily, for all \(x \in \bigoplus_{\iota \in I} \mathbf{E}_{\iota}\) ,

\[
g (x) = \sum_ {i} g (j _ {i} (\mathrm{pr} _ {i} (x))) = \sum_ {i} f _ {i} (\mathrm{pr} _ {i} (x)),
\]

whence the uniqueness of \(g\) . Conversely, setting \(g(x) = \sum_{i} f_i(\mathrm{pr}_i(x))\) for all \(x \in \bigoplus_{i \in I} E_i\) , it is immediately verified that a linear mapping has been defined satisfying the conditions of the statement.

When no confusion can arise, we write \(g = \sum_{i \in I} f_i\) (which is contrary to the conventions of I, §2, no. 1, when the family \((f_i)\) is not of finite support).

In particular, if J is any subset of I, the canonical injections \(j_{i}\) for \(i \in J\) define a canonical linear mapping \(j_{J}: \bigoplus_{i \in J} E_{i} \to \bigoplus_{i \in I} E_{i}\) , which associates with each \((x_{i})_{i \in I}\) the element \((x_{i}')_{i \in I}\) such that \(x_{i}' = x_{i}\) for \(i \in J\) , \(x_{i}' = 0\) for \(i \notin J\) ; this mapping is obviously injective. Moreover, if \((J_{\lambda})_{\lambda \in L}\) is a partition of I, the mapping \(i: \bigoplus_{\lambda \in L} \left( \bigoplus_{i \in J_{\lambda}} E_{i} \right) \to \bigoplus_{i \in I} E_{i}\) corresponding to the family \((j_{J_{\lambda}})\) by Proposition 6 is an isomorphism called canonical (``associativity'' of direct sums).

COROLLARY 1. Let \((\mathbf{E}_{\iota})_{\iota \in I}, (\mathbf{F}_{\lambda})_{\lambda \in L}\) be two families of A-modules. The mapping

\[
\operatorname{Hom} _ {\mathbf {A}} \left(\bigoplus_ {t \in I} \mathrm{E} _ {t}, \prod_ {\lambda \in L} \mathrm{F} _ {\lambda}\right)\rightarrow \prod_ {(t, \lambda) \in I \times L} \operatorname{Hom} _ {\mathbf {A}} \left(\mathrm{E} _ {t}, \mathrm{F} _ {\lambda}\right)\tag{22}
\]

which associates with each \(g \in \mathrm{Hom}_{\mathbf{A}} \left( \bigoplus_{t \in I} \mathbf{E}_t, \prod_{\lambda \in L} \mathbf{F}_\lambda \right)\) the family \((\mathbf{pr}_\lambda \circ g \circ j_t)\) , is a Z-module isomorphism (called canonical).

This follows from Proposition 6 and no. 5, Proposition 4.

COROLLARY 2. Let \((\mathbf{E}_t)_{t \in I}\) be a family of A-modules, \(F\) an A-module and, for each \(t \in I\) , let \(f_t: \mathbf{E}_t \to F\) be a linear mapping. For \(f = \sum_{t \in I} f_t\) to be an isomorphism of \(\mathbf{E} = \bigoplus_{t \in I} \mathbf{E}_t\) onto \(F\) , it is necessary and sufficient that there exist for each \(t \in I\) a linear mapping \(g_t: \mathbf{F} \to \mathbf{E}_t\) with the following properties:

\begin{enumerate}
\def\labelenumi{(\arabic{enumi})}
\item
  \(g_{i} \circ f_{i} = 1_{E_{i}}\) for all \(i \in I\) .
\item
  \(g_{\iota} \circ f_{\kappa} = 0\) for \(\iota \neq \kappa\) .
\item
  For all \(y \in \mathbf{F}\) , the family \((g_{\iota}(y))\) has finite support and
\end{enumerate}

\[
y = \sum_ {i \in I} f _ {i} (g _ {i} (y)).\tag{23}
\]

Note that if I is finite, the last condition may also be written as

\[
\sum_ {i \in I} f _ {i} \circ g _ {i} = 1 _ {\mathbf {F}}.\tag{24}
\]

Obviously the conditions are necessary for they are satisfied by the \(g_{\iota} = \mathrm{pr}_{\iota} \circ f\) . Conversely if they hold, for all \(y \in \mathbf{F}\) , \(g(y) = \sum_{\iota} j_{\iota}(g_{\iota}(y))\) is defined and it is immediate that \(g\) is a linear mapping of \(\mathbf{F}\) into \(\mathbf{E}\) . For all \(y \in \mathbf{F}\) , \(f(g(y)) = \sum_{\iota \in \mathbf{I}} f_{\iota}(g_{\iota}(y)) = y\) by hypothesis. On the other hand, for all \(x \in \mathbf{E}\) , \(g_{\kappa}(f(x)) = g_{\kappa}\left(\sum_{\iota} f_{\iota}(\mathrm{pr}_{\iota}(x))\right) = g_{\kappa}(f_{\kappa}(\mathrm{pr}_{\kappa}(x))) = \mathrm{pr}_{\kappa}(x)\) by hypothesis;

therefore \(g(f(x)) = \sum_{i} j_i(g_i(f(x))) = \sum_{i} j_i(\mathrm{pr}_i(x)) = x\) , which proves the corollary.

PROPOSITION 7. (i) Let \((\mathbf{E}_t)_{t \in I}\) , \((\mathbf{F}_t)_{t \in I}\) be two families of A-modules with the same indexing set I; for every family of linear mappings \(f_t: \mathbf{E}_t \to \mathbf{F}_t\) ( \(t \in I\) ), the restriction to \(\bigoplus_{t \in I} \mathbf{E}_t\) of the linear mapping \((x_t) \mapsto (f_t(x_t))\) is a linear mapping \(f: \bigoplus_{t \in I} \mathbf{E}_t \to \bigoplus_{t \in I} \mathbf{F}_t\) denoted by \(\bigoplus_{t \in I} f_t\) or \(\bigoplus_{t} f_t\) (where \(f = f_p \oplus f_{p+1} \oplus \cdots \oplus f_q\) if \(\mathbf{I} = [p, q]\) is an interval in \(\mathbf{Z}\) ).

\begin{enumerate}
\def\labelenumi{(\roman{enumi})}
\setcounter{enumi}{1}
\tightlist
\item
  Let \((\mathbf{G}_{\iota})_{\iota \in \mathbf{I}}\) be a third family of A-modules with \(\mathbf{I}\) as indexing set and, for all \(\iota \in \mathbf{I}\) , let \(g_{\iota}: F_{\iota} \to G_{\iota}\) be a linear mapping; we write \(g = \bigoplus_{\iota} g_{\iota}\) . For each of the sequences \(\mathbf{E}_{\iota} \xrightarrow{f_{\iota}} F_{\iota} \xrightarrow{g_{\iota}} G_{\iota}\) to be exact, it is necessary and sufficient that the sequence
\end{enumerate}

\[
\bigoplus_ {i} \mathrm{E} _ {i} \xrightarrow {f} \bigoplus_ {i} \mathrm{F} _ {i} \xrightarrow {g} \bigoplus_ {i} \mathrm{G} _ {i}
\]

be exact.

Obviously, for all \((x_{i})\in \bigoplus_{i}\mathbf{E}_{i}\) , the family \((f_{i}(x_{i}))\) has finite support, whence (i). On the other hand, to say that an element \(y = (y_{i})\) of \(\bigoplus_{i}\mathbf{F}_{i}\) belongs to \(\operatorname {Ker}(g)\) means that \(y_{i}\in \operatorname {Ker}(g_{i})\) for all \(i \in I\) (no. 5, Proposition 5); similarly, if \(y_{i}\in \operatorname {Im}(f_{i})\) for all \(\iota \in I\) , there exists for each \(\iota \in I\) an \(x_{i}\in \mathbf{E}_{i}\) such that \(y_{i} = f_{i}(x_{i})\) and when \(y_{i} = 0\) , it may be supposed that \(x_{i} = 0\) ; hence \(y\in \operatorname {Im}(f)\) and the converse is obvious.

COROLLARY 1. Under the conditions of Proposition 7, (i),

\[
\operatorname{Ker} (f) = \bigoplus_ {i \in I} \operatorname{Ker} (f _ {i}), \quad \operatorname{Im} (f) = \bigoplus_ {i \in I} \operatorname{Im} (f _ {i})\tag{25}
\]

and there are canonical isomorphisms

\[
\operatorname{Coim} (f) \cong \bigoplus_ {i \in I} \operatorname{Coim} (f _ {i}), \quad \operatorname{Coker} (f) \cong \bigoplus_ {i \in I} \operatorname{Coker} (f _ {i})
\]

defined as in no. 5, Corollary to Proposition 5. In particular, for \(f\) to be injective (resp. surjective, bijective, zero), it is necessary and sufficient that each of the \(f_i\) be injective (resp. surjective, bijective, zero).

If, for all \(t \in I\) , we consider a submodule \(F_t\) of \(E_t\) , the module \(\bigoplus_{t \in I} F_t\) is a submodule of \(\bigoplus_{t \in I} E_t\) and, by virtue of Corollary 1 to Proposition 7, there is a canonical isomorphism

\[
\bigoplus_ {i \in I} \left(\mathrm{E} _ {i} / \mathrm{F} _ {i}\right)\cong \left(\bigoplus_ {i \in I} \mathrm{E} _ {i}\right) / \left(\bigoplus_ {i \in I} \mathrm{F} _ {i}\right).\tag{26}
\]

COROLLARY 2. Let \((\mathbf{E}_t)_{i \in I}, (\mathbf{E}_t')_{i \in I}, (\mathbf{F}_\lambda)_{\lambda \in L}, (\mathbf{F}_\lambda')_{\lambda \in L}\) be four families of A-modules and, for each \(i \in I\) (resp. each \(\lambda \in L\) ), \(u_i: \mathbf{E}_i' \to \mathbf{E}_i\) (resp. \(v_\lambda: \mathbf{F}_\lambda \to \mathbf{F}_\lambda'\) ) a linear mapping. Then the diagram

\[
\begin{array}{c} \operatorname{Hom} \left(\bigoplus_ {t \in I} E _ {t} ^ {\prime}, \prod_ {\lambda \in L} F _ {\lambda} ^ {\prime}\right) \xrightarrow {\phi^ {\prime}} \prod_ {(t, \lambda) \in I \times L} \operatorname{Hom} (E _ {t} ^ {\prime}, F _ {\lambda} ^ {\prime}) \\ \operatorname{Hom} \left(\bigoplus_ {t} u _ {t}, \prod_ {\lambda} v _ {\lambda}\right) \Bigg | ^ {\uparrow} \\ \operatorname{Hom} \left(\bigoplus_ {t \in I} E, \prod_ {\lambda \in L} F _ {\lambda}\right) \xrightarrow [ \phi ]{\quad} \prod_ {(t, \lambda) \in I \times L} \operatorname{Hom} (E _ {t}, F _ {\lambda}) \end{array}
\]

(where \(\phi\) and \(\phi'\) are the canonical isomorphisms defined in Corollary 1 to Proposition 6) is commutative.

The verification follows immediately from the definitions.

When all the \(E_{t}\) are equal to the same A-module E, the direct sum \(\bigoplus_{i\in I}E_{i}\) is also denoted by \(\mathrm{E}^{(I)}\) : its elements are the mappings of I into E with finite support. If, for all \(t,f_{t}\) is taken to be the identity mapping \(E\to E\) , by Proposition 6, a linear mapping \(\mathrm{E}^{(I)}\to\mathrm{E}\) is obtained, called codiagonal, which associates with every family \((x_{t})_{t\in I}\) of elements of E, of finite support, its sum \(\sum_{i\in I}x_{i}\) .

Remark. Recall that the definition of direct sum extends immediately to a family \((\mathrm{E}_{\iota})_{\iota\in\mathbb{I}}\) of groups which are not necessarily commutative, multiplicative notation of course replacing the additive notation; we then say ``restricted product'' or ``restricted sum'' instead of ``direct sum'' (I, §4, no. 9).

Note that E is a normal subgroup of the product \(F = \prod_{i \in I} E_{i}\) and that each of the \(j_{\kappa}(E_{\kappa})\) is a normal subgroup of F; moreover, for two distinct indices \(\lambda, \mu\) , every element of \(j_{\lambda}(E_{\lambda})\) is permutable with every element of \(j_{\mu}(E_{\mu})\) . Proposition 6 extends to the general case with the hypothesis that, for two distinct indices \(\lambda, \mu\) , every element of \(f_{\lambda}(E_{\lambda})\) is permutable in M with every element of \(f_{\mu}(E_{\mu})\) (I, § 4, no. 9, Proposition 12). The property of ``associativity'' of the restricted sum follows immediately from this. Proposition 7 and its Corollaries 1 and 2 hold without modification.

